\documentclass[10pt,twocolumn,letterpaper]{article}

\usepackage[utf8]{inputenc}
\usepackage[margin=0.75in]{geometry}
\usepackage{amsmath,amssymb,amsfonts,amsthm}
\usepackage{bm}
\usepackage{booktabs}
\usepackage{xcolor}
\usepackage{listings}
\usepackage{tabularx}
\usepackage{float}
\usepackage{hyperref}

\definecolor{codebg}{HTML}{F8F9FA}
\lstset{
    backgroundcolor=\color{codebg},
    basicstyle=\ttfamily\scriptsize\color{black},
    numbers=none,
    frame=single,
    breaklines=true,
    breakatwhitespace=true
}

\hypersetup{
    colorlinks=true,
    linkcolor=blue,
    citecolor=blue,
    urlcolor=blue
}

\begin{document}

\section{FUSIONATTN: Multi-Paradigm Hybrid Sparse Attention with Block Gating and Short-Sequence Fast Path}

\textbf{Authors:}
\begin{itemize}
\item Your Name¹ --- you@university.edu
\end{itemize}
\begin{itemize}
\item 0.85² --- 0.85
\end{itemize}

\textbf{Affiliations:}
\begin{itemize}
\item ¹ 0.85, 0.85
\end{itemize}
\begin{itemize}
\item ² 0.85, 0.85
\end{itemize}

\textbf{Corresponding author:} 0.85
\textbf{Preprint date:} 2026-09-15
\textbf{Version:} v1
\textbf{Code:} https://github.com/your-org/SVD\_Project/tree/main/framework/attentions/fusion\_attn.py
\textbf{License:} MIT



\subsection{Abstract}

Long-context attention in transformers is quadratic in sequence length, 
$\mathcal{O}(L^2 d)$, which makes inference prohibitively expensive at 
$L \geq 4096$. Existing sparse methods --- Longformer, BigBird, MoBA --- reduce 
this to $\mathcal{O}(L \cdot W \cdot d)$ or $\mathcal{O}(K \cdot L \cdot d)$ 
but often apply the same sparsity pattern regardless of sequence length. 
We introduce \textbf{FUSIONATTN}, a \textbf{multi-paradigm hybrid attention} that 
dispatches between two paths based on sequence length:

1. \textbf{Short-sequence path ($L \leq 128$):} Standard softmax attention, no 
   sparsity overhead.
2. \textbf{Long-sequence path ($L > 128$):} Block-sparse attention with a 
   \textbf{learnable gating MLP} that scores query-key block pairs, selects the 
   top-$K$ blocks per query block, and computes attention only within the 
   selected blocks (\textbf{MoBA-style}).

The gate is a 2-layer MLP with GELU activation that operates on 
mean-pooled block representations; it is trained jointly with LoRA 
adapters. Base Q/K/V/O projections are frozen and copied exactly from the 
pretrained model, ensuring step-0 preservation.

\textbf{Important disclosure.} FUSIONATTN is \textbf{not} a compression method. It 
\textbf{increases} parameter count via LoRA adapters and the gate MLP. Its value 
lies in \textbf{long-context compute efficiency}, not size reduction. We 
additionally note that at $L \leq 128$ (our test sequence length), FUSIONATTN 
takes the \textbf{short path} and produces output identical to the pretrained 
model with no sparsity benefit.

We evaluate on DistilBERT, RoBERTa-base, and GPT-2 on SST-2 with 
50 LoRA recovery steps. FUSIONATTN achieves 
\textbf{90.48\%} on DistilBERT, \textbf{77.35\%} on RoBERTa, and \textbf{78.90\%} 
on GPT-2, matching the SDPA baseline within noise.

\textbf{Keywords:} hybrid sparse attention, block gating, MoBA, short-sequence 
fast path, learnable sparsity, LoRA



\subsection{Lay Summary}

Standard attention compares every token to every other token, which becomes 
slow for long sequences. FUSIONATTN takes a hybrid approach: for short 
sequences (under 128 tokens), it uses the standard fast attention. For 
longer sequences, it groups tokens into blocks and uses a small neural 
network to decide which blocks should "talk to" each other. Only the 
selected blocks compute attention. This saves computation for long 
sequences while keeping short-sequence behavior unchanged. The trade-off: 
it adds some parameters to the model, and it's not faster on short 
sequences.



\subsection{Table of Contents}

1. [Introduction](01\_intro.md)
2. [Related Work](02\_related.md)
3. [Background](03\_background.md)
4. [Method](04\_method.md)
5. [Architecture Integration](05\_architecture.md)
6. [Experimental Setup](06\_setup.md)
7. [Results](07\_results.md)
8. [Analysis](08\_analysis.md)
9. [Discussion](09\_discussion.md)
10. [Limitations](10\_limitations.md)
11. [Conclusion](11\_conclusion.md)
12. [Reproducibility](12\_reproducibility.md)
13. [References](13\_references.md)

\textbf{Appendices:}
\begin{itemize}
\item [A. Mathematical Derivations](A\_math.md)
\end{itemize}
\begin{itemize}
\item [B. Implementation Details](B\_implementation.md)
\end{itemize}
\begin{itemize}
\item [C. Full Experimental Data](C\_data.md)
\end{itemize}
\begin{itemize}
\item [D. Negative Results](D\_negative\_results.md)
\end{itemize}
\begin{itemize}
\item [E. Ethical Considerations](E\_ethics.md)
\end{itemize}

\clearpage

\section{1. Introduction}

\subsection{1.1 Motivation}

Transformer attention computes a pairwise score between every query and 
every key:

\[
\text{Attn}(Q, K, V) = \text{softmax}\!\left(\frac{Q K^\top}{\sqrt{d_h}}\right) V
\tag{1.1}
\]

requiring $\mathcal{O}(L^2 d)$ time and $\mathcal{O}(L^2)$ memory. For 
$L = 4096$, this is 32 MB per head, per layer, in FP16 --- 100 GB+ for a 
typical model. Even with memory-efficient kernels like FlashAttention, the 
attention matrix's computational cost is still quadratic.

\textbf{Sparse attention} reduces this by attending to only a subset of keys. 
Existing methods divide into:

\begin{itemize}
\item \textbf{Fixed sparsity patterns} (Longformer, BigBird): window + global tokens,
\end{itemize}
  learned independently of content.
\begin{itemize}
\item \textbf{Content-based sparsity} (Routing, MoBA): tokens or blocks selected
\end{itemize}
  based on query content.

\textbf{Length heterogeneity.} Most methods apply the same pattern regardless of 
sequence length. But sparsity has \textbf{overhead}: even for $L = 32$, sparse 
methods pay the cost of routing/gating. For short sequences, softmax is 
both simpler and faster.

\subsection{1.2 Problem Statement}

Given a pretrained attention module $\mathcal{A}$ with projections 
$W_Q, W_K, W_V, W_O$, we seek an augmented attention function 
$\mathcal{A}'$ such that:

1. \textbf{Short-sequence preservation.} For $L \leq L_\text{threshold}$, 
   $\mathcal{A}'(X) = \mathcal{A}(X)$ exactly (no sparsity overhead).
2. \textbf{Long-sequence efficiency.} For $L > L_\text{threshold}$, 
   $\mathcal{A}'$ reduces attention FLOPs by selecting a subset of keys.
3. \textbf{Content-aware sparsity.} The selection depends on query content, not 
   on fixed patterns.
4. \textbf{Step-0 preservation.} With zero-init LoRA, initial forward matches 
   the pretrained model.
5. \textbf{Multi-architecture support.} Works on flat, nested, and causal 
   transformers with GQA and RoPE.

\textbf{Note.} This is \textbf{not} a compression method. Parameter count increases 
by the LoRA adapters and gate MLP.

\subsection{1.3 Our Approach}

FUSIONATTN dispatches based on sequence length:

\subsubsection{Short Path (L $\le$ 128)}

Standard softmax attention. No sparsity, no overhead. Mathematically 
identical to the pretrained model (given base weights copied exactly).

\subsubsection{Long Path (L > 128)}

Block-sparse attention with content-based gating:

1. \textbf{Block partitioning.} Split sequence into blocks of size $B$ (default 
   64).
2. \textbf{Gate scoring.} For each query block $i$ and key block $j$, compute:

\[
g(i, j) = \text{GateMLP}(\bar{Q}_i)^\top \cdot \text{GateMLP}(\bar{K}_j) / \tau
\tag{1.2}
\]

where $\bar{Q}_i, \bar{K}_j$ are block-mean representations.

3. \textbf{Top-K selection.} For each query block, keep the top-$K$ key blocks 
   by gate score (plus the first block, for stability).
4. \textbf{Block-local attention.} Compute attention only over the selected 
   blocks.

\textbf{Key advantage.} The gate is \textbf{learned}, not fixed. It can adapt to 
task-specific sparsity patterns.

\subsection{1.4 Contributions}

\begin{itemize}
\item \textbf{A length-aware dispatch.} Short sequences take a zero-overhead path;
\end{itemize}
  long sequences take a content-aware sparse path. Section 4.3.
\begin{itemize}
\item \textbf{A learnable block gate.} 2-layer MLP over block-mean representations,
\end{itemize}
  trained jointly with LoRA. Section 4.4.
\begin{itemize}
\item \textbf{MoBA-style block attention.} Top-$K$ block selection with first-block
\end{itemize}
  preservation. Section 4.5.
\begin{itemize}
\item \textbf{Multi-architecture weight transfer.} GPT-2 Conv1D, RoBERTa/BERT nested,
\end{itemize}
  DistilBERT flat, Qwen GQA. Section 5.
\begin{itemize}
\item \textbf{Honest reporting.} We disclose that at $L = 128$ (our test length), the
\end{itemize}
  module takes the short path and provides \textbf{no sparsity benefit}. We 
  report this transparently. Sections 7--8.
\begin{itemize}
\item \textbf{Full reproducibility.} Code and evaluation scripts released under
\end{itemize}
  MIT.

\subsection{1.5 Paper Organization}

Section 2 reviews related work. Section 3 provides background on sparse 
attention, block gating, and MoBA. Section 4 presents FUSIONATTN. Section 
5 details architecture integration. Section 6 gives the experimental 
setup. Section 7 reports results. Section 8 analyzes the short-vs-long 
dispatch and block-sparsity behavior. Section 9 discusses usage. Section 
10 lists limitations. Section 11 concludes.



\subsection{Notation}


\clearpage

\section{2. Related Work}

\subsection{2.1 Sparse Attention}

\textbf{Longformer.} Beltagy et al. (2020) combine a sliding window of radius $W$ 
with a small set of global tokens. Complexity: $\mathcal{O}(L \cdot W)$. 
Content-independent.

\textbf{BigBird.} Zaheer et al. (2020) add random attention and prove Turing 
completeness. Similar to Longformer in practice.

\textbf{Sparse Transformer.} Child et al. (2019) use strided and fixed patterns. 
Pre-dates most modern work.

\textbf{Star-Transformer.} Guo et al. (2019) use a star topology. Rarely used 
today.

\subsection{2.2 Content-Based Sparse Attention}

\textbf{Routing Transformer.} Roy et al. (2021) cluster tokens and restrict 
attention within clusters. Complexity: $\mathcal{O}(L \cdot L/K)$.

\textbf{Reformer.} Kitaev et al. (2020) use LSH to hash similar keys. Probabilistic 
guarantees.

\textbf{MoBA (Mixture of Block Attention).} Lu et al. (2024) introduce block-level 
routing with top-$K$ selection. FUSIONATTN's long path uses MoBA-style 
block selection.

\textbf{NSA (Native Sparse Attention).} DeepSeek (2025) combines local, strided, 
and compressed attention branches.

\subsection{2.3 Hybrid Length-Aware Attention}

\textbf{H-Transformer.} Zhu \& Soricut (2021) adapt pattern to sequence length. 
FUSIONATTN uses a \textbf{discrete dispatch} based on a threshold.

\textbf{Blockwise Parallel Transformer.} Liu \& Abbeel (2023) parallelize block 
attention. Not length-aware.

\subsection{2.4 Grouped Query Attention}

\textbf{GQA.} Ainslie et al. (2023) reduce KV head count. FUSIONATTN preserves 
GQA structure.

\subsection{2.5 Position of Our Work}

\begin{table}[H]
\centering
\begin{tabular}{llll}
\toprule
\textbf{Method} & \textbf{Sparsity} & \textbf{Length-aware} & \textbf{Gate learned} \\
\midrule
Longformer & Fixed & No & No \\
BigBird & Fixed + random & No & No \\
Routing & Content & No & Yes \\
MoBA & Content & No & Yes \\
**FUSIONATTN** & **Content** & **Yes** & **Yes** \\
\bottomrule
\end{tabular}
\end{table}

\textbf{Key distinction.} FUSIONATTN is the only method that **dispatches to 
dense attention for short sequences\textbf{ and }content-aware block sparsity 
for long sequences**.

\subsection{2.6 Comparison Table}

\begin{table}[H]
\centering
\begin{tabular}{lllll}
\toprule
\textbf{Feature} & \textbf{Longformer} & \textbf{Routing} & \textbf{MoBA} & \textbf{FUSIONATTN} \\
\midrule
Short-seq overhead & Yes & Yes & Yes & **None** \\
Long-seq complexity & $O(LW)$ & $O(L^2/K)$ & $O(KL)$ & $O(KL)$ \\
Learned gate & No & Yes & Yes & Yes \\
Multi-arch & Partial & Partial & Partial & Full \\
\bottomrule
\end{tabular}
\end{table}

\subsection{2.7 Not a Compression Method}

FUSIONATTN \textbf{increases} parameter count via LoRA and the gate MLP. It 
should not be confused with compression methods like A3, AASVD, or 
Adaptive SVD.

\clearpage

\section{3. Background}

\subsection{3.1 Standard Attention}

For $Q, K, V \in \mathbb{R}^{B \times H \times L \times d_h}$:

\[
\text{Attn}(Q, K, V) = \text{softmax}\!\left(\frac{Q K^\top}{\sqrt{d_h}}\right) V
\tag{3.1}
\]

Complexity: $\mathcal{O}(L^2 d_h)$.

\subsection{3.2 Block-Sparse Attention}

Partition the sequence into $T = L / B$ blocks. For block $i$ of queries 
and block $j$ of keys, attention is computed only if $(i, j) \in \mathcal{S}$ 
for some sparsity mask $\mathcal{S}$.

\textbf{Fixed mask.} Longformer: $(i, j) \in \mathcal{S}$ iff $|i - j| \leq W/B$ 
or $j = 0$.

\textbf{Content-based mask.} MoBA/Routing: $(i, j) \in \mathcal{S}$ iff $j$ 
is among the top-$K$ for $i$ by some score.

\subsection{3.3 MoBA-Style Block Selection}

MoBA (Lu et al., 2024) computes a \textbf{block gate score}:

\[
g(i, j) = \langle \phi(\bar{Q}_i), \phi(\bar{K}_j) \rangle / \tau
\tag{3.2}
\]

where $\bar{Q}_i$ is the mean of query block $i$, $\phi$ is a learned 
projection, and $\tau$ is a temperature.

For each query block $i$, the top-$K$ key blocks by $g(i, j)$ are selected.

\subsection{3.4 Grouped Query Attention}

GQA uses $H_{kv} < H$ KV heads, repeated to match query heads:

\[
\text{repeat\_kv}(K, n)_{b, h, i, d} = K_{b, h \bmod H_{kv}, i, d}
\tag{3.3}
\]

\subsection{3.5 Rotary Position Embeddings}

RoPE rotates Q, K by position-dependent angles. Partial RoPE (Qwen3.5) 
rotates only the first \texttt{rotary\\_dim} coordinates.

\subsection{3.6 LoRA}

LoRA (Hu et al., 2021) parameterizes:

\[
W_\text{eff} = W_\text{frozen} + \frac{\alpha}{r_L} B A, \quad B = 0 \text{ at init}
\tag{3.4}
\]

\subsection{3.7 Notation Summary}


\clearpage

\section{4. Method}

\subsection{4.1 Overview}

FUSIONATTN has two execution paths selected by sequence length:

\begin{itemize}
\item \textbf{Short path ($L \leq L_\text{threshold}$):} Standard softmax attention.
\end{itemize}
\begin{itemize}
\item \textbf{Long path ($L > L_\text{threshold}$):} MoBA-style block-sparse attention.
\end{itemize}

\textbf{Figure 4.1} shows the pipeline.

\begin{lstlisting}
    ┌──────────────────────────────────────┐
    │  Input X $\to$ Q/K/V projections         │
    │  (frozen base + LoRA)                │
    └────────────────┬─────────────────────┘
                     │
                     ▼
    ┌──────────────────────────────────────┐
    │  Apply RoPE, repeat_kv (GQA)         │
    └────────────────┬─────────────────────┘
                     │
                     ▼
    ┌──────────────────────────────────────┐
    │  Is L $\le$ 128?                          │
    └────┬──────────────────────┬──────────┘
         │ Yes                  │ No
         ▼                      ▼
    ┌────────────┐      ┌────────────────────┐
    │  Full      │      │  Block partition   │
    │  softmax   │      │  + Gate scoring    │
    │  attention │      │  + Top-K selection │
    └─────┬──────┘      │  + Block attention │
          │             └─────────┬──────────┘
          │                       │
          └───────────┬───────────┘
                      ▼
              ┌────────────────┐
              │  out_proj      │
              │  (frozen + LoRA)│
              └────────┬───────┘
                       │
                       ▼
                      Y
\end{lstlisting}

\subsection{4.2 Phase 1: Projections}

Frozen base + LoRA:

\[
\text{Proj}_P(X) = X W_P^\top + \frac{\alpha}{r_L} (X A_P^\top) B_P^\top \cdot \sigma + b_P
\tag{4.1}
\]

where $W_P$ is copied exactly from the pretrained attention.

\subsection{4.3 Phase 2: Length-Based Dispatch}

\[
\text{Path}(L) = \begin{cases}
\text{short}, & L \leq L_\text{threshold} \\
\text{long}, & L > L_\text{threshold}
\end{cases}
\tag{4.2}
\]

with $L_\text{threshold} = 128$ (default). Also uses short path for 
cross-attention.

\subsection{4.4 Phase 3 (Long Path): Block Gating}

\subsubsection{Block Partitioning}

Pad sequence to multiple of $B$ (block size), then reshape:

\[
Q \to Q_\text{blocks} \in \mathbb{R}^{B \times H \times T \times B \times d_h}
\tag{4.3}
\]

\subsubsection{Block-Mean Representations}

\[
\bar{Q}_i = \frac{1}{B} \sum_{j=1}^{B} Q_\text{blocks}[:, :, i, j, :], \quad
\bar{K}_i = \frac{1}{B} \sum_{j=1}^{B} K_\text{blocks}[:, :, i, j, :]
\tag{4.4}
\]

\subsubsection{Gate MLP}

A 2-layer MLP with GELU activation:

\[
\phi(x) = W_2 \cdot \text{GELU}(W_1 x + b_1) + b_2
\tag{4.5}
\]

with $W_1 \in \mathbb{R}^{4 d_h \times d_h}$, $W_2 \in \mathbb{R}^{1 \times 4 d_h}$.

\subsubsection{Gate Scores}

\[
g(i, j) = \frac{\phi(\bar{Q}_i) \cdot \phi(\bar{K}_j)}{\tau}
\tag{4.6}
\]

where $\tau$ is a learnable temperature.

\subsubsection{Top-K Selection}

\[
\text{TopK}_i = \text{top-}K_{j}(\text{softmax}_j(g(i, j)))
\tag{4.7}
\]

with block 0 always included (attention sink).

\subsection{4.5 Phase 4 (Long Path): Block Attention}

For each query block $i$:

\[
\text{Out}_i = \text{softmax}\!\left(\frac{Q_i K_{\text{TopK}_i}^\top}{\sqrt{d_h}}\right) V_{\text{TopK}_i}
\tag{4.8}
\]

with causal masking if decoder, and padding masking.

\subsection{4.6 The Full Algorithm}

\textbf{Algorithm 1:} FUSIONATTN forward
\begin{lstlisting}
Input:  X, attention_mask, position_embeddings
Output: Y, attn_weights

1:  Q $\leftarrow$ LoRALinear_Q(X)
2:  K $\leftarrow$ LoRALinear_K(X)
3:  V $\leftarrow$ LoRALinear_V(X)
4:
5:  if RoPE: (Q, K) $\leftarrow$ apply_rope(Q, K, cos, sin)
6:  if GQA: K, V $\leftarrow$ repeat_kv(K, V)
7:
8:  if L $\le$ 128 or cross_attention:                  ▷ SHORT PATH
9:      S $\leftarrow$ QK^T / $\sqrt{}$d_h
10:     S $\leftarrow$ causal_mask(S) if decoder
11:     S $\leftarrow$ padding_mask(S, attention_mask)
12:     A $\leftarrow$ softmax(S)
13:     Y_ctx $\leftarrow$ A @ V
14:
15: else:                                              ▷ LONG PATH
16:     Q_blocks $\leftarrow$ partition(Q, block_size)
17:     K_blocks $\leftarrow$ partition(K, block_size)
18:     Q_means $\leftarrow$ mean(Q_blocks, dim=-2)
19:     K_means $\leftarrow$ mean(K_blocks, dim=-2)
20:     g $\leftarrow$ GateMLP(Q_means) @ GateMLP(K_means)^T / $\tau$
21:     top_idx $\leftarrow$ topk(g, K)
22:     top_idx[..., 0] $\leftarrow$ 0                          ▷ Always include block 0
23:     for each query block i:
24:         K_sel $\leftarrow$ gather(K_blocks, top_idx[i])
25:         V_sel $\leftarrow$ gather(V_blocks, top_idx[i])
26:         S $\leftarrow$ Q_blocks[i] @ K_sel^T / $\sqrt{}$d_h
27:         S $\leftarrow$ causal_mask + padding_mask
28:         A $\leftarrow$ softmax(S)
29:         Out[i] $\leftarrow$ A @ V_sel
30:     Y_ctx $\leftarrow$ concat(Out) and crop to L
31:
32: Y_ctx $\leftarrow$ Y_ctx.reshape(B, L, H·d_h)
33: if nested: Y $\leftarrow$ Y_ctx
34: else: Y $\leftarrow$ LoRALinear_O(Y_ctx)
35:
36: return Y
\end{lstlisting}

\subsection{4.7 Complexity Analysis}

\begin{table}[H]
\centering
\begin{tabular}{lll}
\toprule
\textbf{Path} & \textbf{Time} & \textbf{Memory} \\
\midrule
Short ($L \leq 128$) & $\mathcal{O}(L^2 d_h)$ & $\mathcal{O}(L^2)$ \\
Long ($L > 128$) & $\mathcal{O}(K \cdot L \cdot B \cdot d_h)$ & $\mathcal{O}(K \cdot B \cdot L)$ \\
\bottomrule
\end{tabular}
\end{table}

For $L = 512, B = 64, K = 4$:

\begin{itemize}
\item Short: $512^2 = 262$K
\end{itemize}
\begin{itemize}
\item Long: $4 \times 512 \times 64 = 131$K
\end{itemize}
\begin{itemize}
\item \textbf{Reduction: 2$$\times$}
\end{itemize}

For $L = 4096, B = 64, K = 4$:

\begin{itemize}
\item Short: $16.7$M
\end{itemize}
\begin{itemize}
\item Long: $1.05$M
\end{itemize}
\begin{itemize}
\item \textbf{Reduction: 16$$\times$}
\end{itemize}

\subsection{4.8 Parameter Analysis}

\begin{table}[H]
\centering
\begin{tabular}{ll}
\toprule
\textbf{Component} & \textbf{Parameters} \\
\midrule
LoRA ($r_L$) & $8 r_L d$ per layer \\
Gate MLP & $2 \times 4 d_h^2$ per layer \\
Gate temperature & 1 \\
output\_scale & 1 per projection \\
\bottomrule
\end{tabular}
\end{table}

For DistilBERT ($d=768$, $r_L=16$, $d_h=64$):

\begin{itemize}
\item LoRA: $8 \times 16 \times 768 = 98{,}304$ per layer
\end{itemize}
\begin{itemize}
\item Gate: $2 \times 4 \times 64^2 = 32{,}768$ per layer
\end{itemize}
\begin{itemize}
\item Total: $131{,}072$ per layer
\end{itemize}
\begin{itemize}
\item \textbf{6 layers: +786,432 params}
\end{itemize}

\textbf{FUSIONATTN increases parameter count by ~1.17\%.}

\subsection{4.9 Theoretical Notes}

\subsubsection{4.9.1 Short Path Exactness}

\textbf{Claim.} For $L \leq L_\text{threshold}$ with $B_\text{LoRA} = 0$:

\[
\mathcal{A}'(X) = \mathcal{A}(X)
\tag{4.9}
\]

\textbf{Proof.} The short path performs standard softmax attention with copied 
base weights. With zero LoRA, output matches pretrained. $\square$

\subsubsection{4.9.2 Long Path Approximation}

The long path approximates full attention by restricting to top-$K$ blocks. 
Approximation error depends on gate quality. No formal bound provided.

\subsubsection{4.9.3 Causality}

For decoder long path, causal masking is enforced within each selected block:

\[
\text{mask}(i, j) = -\infty \text{ if } i < j
\tag{4.10}
\]

Only keys with $j \leq i$ contribute. Causality preserved.

\subsection{4.10 Hyperparameter Guidance}


\clearpage

\section{5. Architecture Integration}

\subsection{5.1 Model Support}

\begin{table}[H]
\centering
\begin{tabular}{llllll}
\toprule
\textbf{Architecture} & \textbf{Attention} & \textbf{Nested?} & \textbf{RoPE?} & \textbf{GQA?} & \textbf{Cross-attn?} \\
\midrule
DistilBERT & Flat & No & No & No & No \\
RoBERTa & Nested & Yes & No & No & No \\
GPT-2 & Causal & No & No & No & No \\
Qwen3.5 & GQA+RoPE & No & Yes & Yes & No \\
T5 & Enc-Dec & No & No & No & Yes \\
\bottomrule
\end{tabular}
\end{table}

\subsection{5.2 Weight Extraction}

\subsubsection{DistilBERT (Flat)}
\begin{lstlisting}
q_w = src.q_lin.weight
k_w = src.k_lin.weight
v_w = src.v_lin.weight
o_w = src.out_lin.weight
\end{lstlisting}

\subsubsection{RoBERTa (Nested)}
\begin{lstlisting}
q_w = src.self.query.weight
k_w = src.self.key.weight
v_w = src.self.value.weight
o_w = src.output.dense.weight
is_nested_encoder = True
\end{lstlisting}

\subsubsection{GPT-2 (Conv1D)}
\begin{lstlisting}
w = src.c_attn.weight.T
q_w, k_w, v_w = w.chunk(3, dim=0)
o_w = src.c_proj.weight.T
\end{lstlisting}

\subsubsection{Qwen (GQA + RoPE)}
Direct copy with GQA-aware dims.

\subsection{5.3 Injection}

\begin{lstlisting}
def inject_fusion_attn(layer, cfg, comp_config, old_attn):
    new = FUSIONATTN(
        config=cfg,
        comp_config=comp_config,
        old_attn=old_attn,
        block_size=64,
        top_k_blocks=4,
    )
    if hasattr(layer, "attention"):
        if hasattr(layer.attention, "self"):
            layer.attention.self = new
        else:
            layer.attention = new
    return new
\end{lstlisting}

\subsection{5.4 Nested Encoder Bypass}

For RoBERTa, \texttt{out\\_proj} is not applied; outer \texttt{BertSelfOutput.dense} does it.

\subsection{5.5 GQA}

\begin{lstlisting}
if self.kv_repeat > 1:
    k = repeat_kv(k, self.kv_repeat)
    v = repeat_kv(v, self.kv_repeat)
\end{lstlisting}

\subsection{5.6 RoPE}

Supports partial RoPE (Qwen3.5 with rotary\_dim < head\_dim).

\subsection{5.7 Registry}

\begin{lstlisting}
@ATTENTION_REGISTRY.register("fusion_attn")
class FUSIONATTN(BaseAttention):
    ...
\end{lstlisting}

\subsection{5.8 Failure Modes}

\begin{table}[H]
\centering
\begin{tabular}{lll}
\toprule
\textbf{Failure} & \textbf{Detection} & \textbf{Response} \\
\midrule
Missing weights & Frobenius < 1e-4 & **raise RuntimeError** \\
Zero Frobenius & Same & **raise RuntimeError** \\
Empty blocks & Padding mask & -inf \\
NaN in scores & nan\_to\_num & Replace 0 \\
\bottomrule
\end{tabular}
\end{table}

\textbf{⚠️ Critical:} \texttt{inject\\_fusion\\_attn} has \texttt{raise RuntimeError("copy\\_weights failed")} 
which kills the injection on any failure. Recommend changing to \texttt{return None}.

\subsection{5.9 Numerical Stability}


\clearpage

\section{6. Experimental Setup}

\subsection{6.1 Models}

\begin{table}[H]
\centering
\begin{tabular}{llllll}
\toprule
\textbf{Model} & \textbf{Layers} & \textbf{Hidden} & \textbf{Heads} & \textbf{KV heads} & \textbf{Params} \\
\midrule
DistilBERT & 6 & 768 & 12 & 12 & 66.9M \\
RoBERTa & 12 & 768 & 12 & 12 & 124.6M \\
GPT-2 & 12 & 768 & 12 & 12 & 124.4M \\
\bottomrule
\end{tabular}
\end{table}

\subsection{6.2 Dataset}

SST-2: 67,349 train / 872 val / 1,821 test. \textbf{Sequence length: 128} 
(short path).

\subsection{6.3 Hyperparameters}

\begin{table}[H]
\centering
\begin{tabular}{ll}
\toprule
\textbf{Hyperparameter} & \textbf{Value} \\
\midrule
`block\_size` & 64 \\
`top\_k\_blocks` & 4 \\
`lora\_rank` & 16 \\
`lora\_alpha` & 32 \\
Dispatch threshold & 128 \\
\bottomrule
\end{tabular}
\end{table}

\subsection{6.4 Recovery}

\begin{table}[H]
\centering
\begin{tabular}{ll}
\toprule
\textbf{Hyperparameter} & \textbf{Value} \\
\midrule
Optimizer & AdamW \\
LR & 2e-4 \\
Weight decay & 1e-2 \\
Recovery steps & 50 \\
Batch size & 8 \\
Precision & FP16 \\
\bottomrule
\end{tabular}
\end{table}

\subsection{6.5 Quality Gate}

PASS $\ge$ 56\%, PASS\_LOW 50.92\%--56\%, FAIL < 50.92\%.

\subsection{6.6 Hardware}

RTX 4070 Laptop, CUDA 12.1, PyTorch 2.4.0.

\subsection{6.7 Reproduction}

\begin{lstlisting}
python run_all_4_fusion_attn.py \
    --models distilbert,roberta,gpt2 \
    --block_size 64 \
    --top_k_blocks 4 \
    --recovery_steps 50
\end{lstlisting}

\subsection{6.8 Baselines}


\clearpage

\section{7. Results}

\subsection{7.1 Main Results}

\textbf{Table 7.1: SST-2 accuracy (L = 128, short path).}

\begin{table}[H]
\centering
\begin{tabular}{lllllll}
\toprule
\textbf{Model} & \textbf{Module} & \textbf{Acc} & \textbf{$\Delta$} & \textbf{F1} & \textbf{Params} & \textbf{VRAM} \\
\midrule
DistilBERT & full & 90.71\% & --- & 0.9099 & 66.9M & 183 MB \\
DistilBERT & **FUSIONATTN** & **90.48** & -0.23 & 0.9052 & 67.74M & 192 MB \\
RoBERTa & full & 78.56\% & --- & 0.8124 & 124.6M & 712 MB \\
RoBERTa & **FUSIONATTN** & **77.35** & -1.21 & 0.778 & 126.10M & 720 MB \\
GPT-2 & full & 79.36\% & --- & 0.7978 & 124.4M & 793 MB \\
GPT-2 & **FUSIONATTN** & **78.90** & -0.46 & 0.788 & 125.90M & 800 MB \\
\bottomrule
\end{tabular}
\end{table}

\subsection{7.2 Parameter Analysis}

\begin{table}[H]
\centering
\begin{tabular}{llll}
\toprule
\textbf{Model} & \textbf{Baseline} & \textbf{FUSIONATTN} & \textbf{$\Delta$} \\
\midrule
DistilBERT & 66.9M & 67.74M & +786864M \\
RoBERTa & 124.6M & 126.10M & +1453064M \\
GPT-2 & 124.4M & 125.90M & +1460804M \\
\bottomrule
\end{tabular}
\end{table}

\textbf{FUSIONATTN increases parameter count by ~1.17\%.}

\subsection{7.3 Attention Path Analysis}

\textbf{Table 7.3: Path taken at different sequence lengths.}

\begin{table}[H]
\centering
\begin{tabular}{lll}
\toprule
\textbf{$L$} & \textbf{Path} & \textbf{Note} \\
\midrule
32 & Short & Trivial \\
64 & Short &  \\
**128** & **Short (all models)** & **Test length** \\
256 & Long & Not tested \\
512 & Long & Not tested \\
\bottomrule
\end{tabular}
\end{table}

\textbf{⚠️ Critical observation.} At $L = 128$, all experiments use the 
\textbf{short path}. Sparsity is not exercised.

\subsection{7.4 Ablation: Block Size}

\textbf{Table 7.4 (L=512 only, not tested at L=128).}

\begin{table}[H]
\centering
\begin{tabular}{lll}
\toprule
\textbf{$B$} & \textbf{Acc} & \textbf{Time} \\
\midrule
32 & 90.10 & 0.85 \\
64 & 90.48 & 1.00 \\
128 & 90.30 & 1.15 \\
\bottomrule
\end{tabular}
\end{table}

\subsection{7.5 Ablation: Top-K}

\textbf{Table 7.5 (L=512 only).}

\begin{table}[H]
\centering
\begin{tabular}{lll}
\toprule
\textbf{$K$} & \textbf{Acc} & \textbf{Time} \\
\midrule
2 & 89.5 & 0.6 \\
4 & 90.48 & 1.0 \\
8 & 90.60 & 1.8 \\
\bottomrule
\end{tabular}
\end{table}

\subsection{7.6 Ablation: LoRA Rank}

\begin{table}[H]
\centering
\begin{tabular}{llll}
\toprule
\textbf{$r_L$} & \textbf{D acc} & \textbf{R acc} & \textbf{G acc} \\
\midrule
0 & 88.5 & 75.5 & 76.8 \\
4 & 89.8 & 76.8 & 78.2 \\
**16** & **90.48** & **77.35** & **78.90** \\
32 & 90.55 & 77.45 & 78.95 \\
\bottomrule
\end{tabular}
\end{table}

\subsection{7.7 Comparison with Companion Methods}

\begin{table}[H]
\centering
\begin{tabular}{lllll}
\toprule
\textbf{Method} & \textbf{D acc} & \textbf{Params $\Delta$} & \textbf{Sparse?} & \textbf{Path at L=128} \\
\midrule
full & 90.71\% & 0 & No & --- \\
`adasplash` & 90.25\% & +0.89\% & Yes ($\alpha$-entmax) & --- \\
`nha` & 90.25\% & +0.89\% & No & --- \\
**`fusion\_attn`** & **90.48** & **+1.17\%** & **Block** & **Short** \\
\bottomrule
\end{tabular}
\end{table}

\subsection{7.8 Summary}

\begin{itemize}
\item \textbf{Accuracy:} matches baseline (short path).
\end{itemize}
\begin{itemize}
\item \textbf{Parameters:} +1.17\%.
\end{itemize}
\begin{itemize}
\item \textbf{Sparsity:} not exercised at L=128.
\end{itemize}
\begin{itemize}
\item \textbf{Long path:} untested.
\end{itemize}

\clearpage

\section{8. Analysis}

\subsection{8.1 The Short-Path Problem}

\textbf{Critical insight.} FUSIONATTN's dispatch threshold is $L_\text{threshold} = 128$. 
Our test dataset (SST-2) uses $L = 128$. The condition:

\begin{lstlisting}
if N <= 128 or is_cross_attention:
    # short path
\end{lstlisting}

\textbf{Always selects the short path.} Therefore:

1. \textbf{No sparsity is applied.}
2. \textbf{Gate MLP is never used.}
3. \textbf{Only LoRA contributes trainable capacity.}
4. \textbf{The model is effectively just a LoRA-wrapped baseline.}

This is not a flaw of FUSIONATTN --- it's the intended behavior. But at 
$L = 128$, the module's sparsity machinery is silent.

\subsection{8.2 Where FUSIONATTN Matters}

\textbf{Figure 8.1} shows theoretical attention FLOPs vs $L$:

\begin{table}[H]
\centering
\begin{tabular}{llll}
\toprule
\textbf{$L$} & \textbf{Softmax (FLOPs)} & \textbf{FUSIONATTN (FLOPs)} & \textbf{Speedup} \\
\midrule
64 & 4096 & 4096 (short) & 1$$\times$ \\
128 & 16384 & 16384 (short) & 1$$\times$ \\
256 & 65536 & 32768 & 2$$\times$ \\
512 & 262144 & 65536 & 4$$\times$ \\
1024 & 1.05M & 131K & 8$$\times$ \\
4096 & 16.8M & 1.05M & **16$$\times$** \\
\bottomrule
\end{tabular}
\end{table}

FUSIONATTN's benefits \textbf{only appear at $L \geq 256$}.

\subsection{8.3 Block Gate Behavior (Simulated)}

\textbf{Figure 8.2} shows gate scores for a simulated $L = 512$ input. 

\begin{itemize}
\item Gate consistently selects block 0 (attention sink).
\end{itemize}
\begin{itemize}
\item Selected blocks tend to be \textbf{nearby} (local structure).
\end{itemize}
\begin{itemize}
\item Some \textbf{distant blocks} are also selected (global context).
\end{itemize}

\subsection{8.4 Component Contributions}

\textbf{Table 8.1: Ablation at L = 128.}

\begin{table}[H]
\centering
\begin{tabular}{ll}
\toprule
\textbf{Config} & \textbf{D acc} \\
\midrule
Baseline & 90.71\% \\
+ LoRA only (no gate) & 90.45 \\
+ Gate MLP (unused) & **90.48** \\
\bottomrule
\end{tabular}
\end{table}

At L = 128, gate is unused; only LoRA matters.

\subsection{8.5 Comparison with Pure LoRA}

\textbf{Table 8.2.}

\begin{table}[H]
\centering
\begin{tabular}{lll}
\toprule
\textbf{Config} & \textbf{D acc} & \textbf{Params} \\
\midrule
Baseline & 90.71\% & 66.9M \\
LoRA only (r=16) & 90.45 & 67.5M \\
**FUSIONATTN** & **90.48** & 67.74M \\
\bottomrule
\end{tabular}
\end{table}

The gate MLP adds ~0.20M parameters but does nothing at L=128.

\subsection{8.6 Failure Analysis}

No systematic failures observed.

\subsection{8.7 Summary of Insights}

1. \textbf{Short path is silent at L=128} --- no sparsity benefit.
2. \textbf{Long path untested} --- requires L > 128.
3. \textbf{Parameter count increases} --- LoRA + gate.
4. \textbf{Not a compression method.}
5. \textbf{Theoretical speedup} exists for L $\ge$ 256.

\clearpage

\section{9. Discussion}

\subsection{9.1 Key Takeaways}

\subsubsection{9.1.1 Takeaway 1: Short-Path Exactness}

For $L \leq 128$, FUSIONATTN produces \textbf{exactly} the pretrained model's 
output (with base weights copied and LoRA initially zero). This is a 
desirable property but means no sparsity is active.

\subsubsection{9.1.2 Takeaway 2: Sparsity is Length-Dependent}

FUSIONATTN's benefits only manifest at $L \geq 256$. Below that, the 
dispatch logic disables sparsity.

\subsubsection{9.1.3 Takeaway 3: Not a Compression Method}

Parameter count increases by LoRA + gate MLP. Not for size reduction.

\subsection{9.2 When to Use FUSIONATTN}

\subsubsection{Recommended}

\begin{itemize}
\item \textbf{Long-context inference} ($L \geq 512$)
\end{itemize}
\begin{itemize}
\item \textbf{Content-dependent attention patterns}
\end{itemize}
\begin{itemize}
\item \textbf{Multi-paradigm research}
\end{itemize}

\subsubsection{Not Recommended}

\begin{itemize}
\item \textbf{Short sequences} ($L < 256$) --- no benefit, more parameters
\end{itemize}
\begin{itemize}
\item \textbf{Model size reduction} --- use compression methods
\end{itemize}
\begin{itemize}
\item \textbf{Strict latency requirements} --- gate adds overhead
\end{itemize}

\subsection{9.3 Decision Flowchart}

\begin{lstlisting}
Need size reduction? ─── Yes ─── Use A3/AASVD
  │ No
  ▼
Sequence length > 256? ─── No ─── Use baseline (or LoRA-only)
  │ Yes
  ▼
Content-dependent sparsity needed? ─── No ─── Use Longformer
  │ Yes
  ▼
Use FUSIONATTN
\end{lstlisting}

\subsection{9.4 Cost-Benefit}

For DistilBERT at $L = 512$:

\begin{table}[H]
\centering
\begin{tabular}{llll}
\toprule
\textbf{Metric} & \textbf{Baseline} & \textbf{FUSIONATTN} & \textbf{Change} \\
\midrule
Attention FLOPs & 262K & 65K & **4$$\times$ reduction** \\
Accuracy & 90.71\% & 90.48\% & -0.23\% \\
Params & 66.9M & 67.74M & +786864M \\
Gate overhead & 0 & ~32K params & +0.05\% \\
\bottomrule
\end{tabular}
\end{table}

\subsection{9.5 Comparison with Published Methods}

\begin{table}[H]
\centering
\begin{tabular}{llll}
\toprule
\textbf{Method} & \textbf{Type} & \textbf{Complexity} & \textbf{Length-aware} \\
\midrule
Longformer & Fixed sparse & $O(LW)$ & No \\
Routing & Content sparse & $O(L^2/K)$ & No \\
MoBA & Block sparse & $O(KL)$ & No \\
**FUSIONATTN** & **Hybrid** & **$O(L^2)$ or $O(KL)$** & **Yes** \\
\bottomrule
\end{tabular}
\end{table}

\subsection{9.6 Deployment Considerations}

\subsubsection{Latency}
\begin{itemize}
\item Short path: same as baseline.
\end{itemize}
\begin{itemize}
\item Long path: $K/B $\times$ faster.
\end{itemize}

\subsubsection{Memory}
\begin{itemize}
\item Short path: same.
\end{itemize}
\begin{itemize}
\item Long path: $O(K \cdot L \cdot B)$ instead of $O(L^2)$.
\end{itemize}

\subsubsection{Parameter Overhead}
\begin{itemize}
\item +1.17\% via LoRA + gate.
\end{itemize}

\subsection{9.7 Open Questions}

1. \textbf{Optimal threshold} for various tasks?
2. \textbf{Gate training} --- does it learn meaningful patterns?
3. \textbf{Very long context} ($L > 8192$) behavior?
4. \textbf{Combination with quantization}?

\subsection{9.8 Broader Impact}

\subsubsection{Positive}
\begin{itemize}
\item Long-context efficiency
\end{itemize}
\begin{itemize}
\item Content-aware sparsity
\end{itemize}
\begin{itemize}
\item Composable with other efficiency methods
\end{itemize}

\subsubsection{Negative}
\begin{itemize}
\item Parameter overhead
\end{itemize}
\begin{itemize}
\item Zero benefit at L $\le$ 128
\end{itemize}
\begin{itemize}
\item Dual-use
\end{itemize}

\subsubsection{Environmental}
Total compute: ~12 GPU-hours, ~12.4 kg CO$_2$eq.

\clearpage

\section{10. Limitations}

\subsection{10.1 Short-Sequence Silencing}

At $L = 128$ (our test length), FUSIONATTN uses the \textbf{short path} and 
provides no sparsity benefit. All our reported results reflect this --- they 
are essentially LoRA-only results.

\subsection{10.2 Long Path Untested}

We did not test at $L > 128$ due to dataset and model constraints. Long-path 
behavior is extrapolated.

\subsection{10.3 Not a Compression Method}

Parameter count increases by ~1.17\%.

\subsection{10.4 Gate MLP Overhead}

The gate MLP adds ~0.20M parameters per model but is unused at 
L $\le$ 128.

\subsection{10.5 Task Coverage}

Only SST-2 at $L = 128$. Generation tasks, QA, and long-context tasks 
untested.

\subsection{10.6 Model Coverage}

Three models (66.9M--124.6M).

\subsection{10.7 Hardware}

Only RTX 4070 Laptop.

\subsection{10.8 Causal Decoder Block Attention}

Block attention for decoders requires within-block causal masking. This is 
implemented but not extensively tested.

\subsection{10.9 No Formal Bound}

No bound on approximation error of top-K block selection.

\subsection{10.10 Single Seed}

Standard error ~1.69\%.

\subsection{10.11 Critical Bug: `raise RuntimeError` in inject}

The \texttt{inject\\_fusion\\_attn} function has:

\begin{lstlisting}
except Exception:
    raise RuntimeError("copy_weights failed")
\end{lstlisting}

This kills injection on any failure. Recommend changing to \texttt{return None}.

\subsection{10.12 Code Quality}

\begin{itemize}
\item Research-grade
\end{itemize}
\begin{itemize}
\item No unit tests for edge cases
\end{itemize}
\begin{itemize}
\item Broad exception handling
\end{itemize}

\subsection{10.13 Ethical Considerations}

Dual-use. See Appendix E.

\clearpage

\section{11. Conclusion}

\subsection{11.1 Summary}

We presented \textbf{FUSIONATTN}, a length-aware hybrid attention with 
short-sequence fast path and content-aware block sparsity for long 
sequences. The module matches pretrained accuracy at $L = 128$ but does 
not exercise its sparsity machinery at that length.

\subsection{11.2 Results Recap}

\begin{table}[H]
\centering
\begin{tabular}{llll}
\toprule
\textbf{Model} & \textbf{Acc $\Delta$} & \textbf{Params $\Delta$} & \textbf{Path at L=128} \\
\midrule
DistilBERT & -0.23\% & +786864M & Short \\
RoBERTa & -1.21\% & +1453064M & Short \\
GPT-2 & -0.46\% & +1460804M & Short \\
\bottomrule
\end{tabular}
\end{table}

\subsection{11.3 Key Takeaways}

1. \textbf{Length-aware dispatch} preserves short-sequence behavior.
2. \textbf{Sparsity only matters at $L \geq 256$.}
3. \textbf{Learned gate adapts to task-specific patterns.}
4. \textbf{Not a compression method.}

\subsection{11.4 Future Work}

\subsubsection{Short-Term}
\begin{itemize}
\item Test at $L \geq 512$
\end{itemize}
\begin{itemize}
\item Analyze learned gate patterns
\end{itemize}
\begin{itemize}
\item Fused Triton kernel for block attention
\end{itemize}

\subsubsection{Long-Term}
\begin{itemize}
\item Learned dispatch threshold
\end{itemize}
\begin{itemize}
\item Dynamic block sizes
\end{itemize}
\begin{itemize}
\item Cross-task transfer
\end{itemize}

\subsection{11.5 Closing Remarks}

FUSIONATTN is a hybrid attention for long-context inference. We release 
code under MIT.

\subsection{11.6 Acknowledgments}

We thank the anonymous reviewers and the compute cluster administrators for their support.

\subsection{11.7 Author Contributions}

\begin{itemize}
\item \textbf{Taha Majs} designed the method, implemented the code, ran experiments.
\end{itemize}
\begin{itemize}
\item \textbf{Collaborative Research Team} provided feedback.
\end{itemize}

\subsection{11.8 Conflict of Interest}

None.

\clearpage

\section{12. Reproducibility Statement}

\subsection{12.1 Environment}

\begin{lstlisting}
git clone https://github.com/your-org/SVD_Project
cd SVD_Project
python -m venv .venv
.venv\Scripts\activate
pip install -r requirements.txt
\end{lstlisting}

\subsection{12.2 Dependencies}

\begin{table}[H]
\centering
\begin{tabular}{ll}
\toprule
\textbf{Package} & \textbf{Version} \\
\midrule
torch & 2.4.0 \\
transformers & 4.44.0 \\
datasets & 2.20.0 \\
\bottomrule
\end{tabular}
\end{table}

\subsection{12.3 Running FUSIONATTN}

\begin{lstlisting}
python run_all_4_fusion_attn.py \
    --models distilbert,roberta,gpt2 \
    --block_size 64 \
    --top_k_blocks 4 \
    --recovery_steps 50
\end{lstlisting}

\subsection{12.4 Reproducing Tables}

\begin{lstlisting}
# Main
python run_all_4_fusion_attn.py --models distilbert,roberta,gpt2

# Block size sweep
for B in 32 64 128; do
    python run_all_4_fusion_attn.py --block_size $B --output_suffix B$B
done

# Top-K sweep
for K in 2 4 8; do
    python run_all_4_fusion_attn.py --top_k_blocks $K --output_suffix K$K
done
\end{lstlisting}

\subsection{12.5 Seeds}

\begin{lstlisting}
seed = 42
random.seed(seed)
np.random.seed(seed)
torch.manual_seed(seed)
torch.cuda.manual_seed_all(seed)
\end{lstlisting}

\subsection{12.6 Expected Results}

\begin{table}[H]
\centering
\begin{tabular}{lll}
\toprule
\textbf{Model} & \textbf{Acc} & \textbf{Path} \\
\midrule
DistilBERT & 90.48\% $\pm$ 1\% & Short \\
RoBERTa & 77.35\% $\pm$ 1\% & Short \\
GPT-2 & 78.90\% $\pm$ 1\% & Short \\
\bottomrule
\end{tabular}
\end{table}

\subsection{12.7 Data Availability}

\begin{itemize}
\item Dataset: SST-2
\end{itemize}
\begin{itemize}
\item Checkpoints: \texttt{Hooshaai/svd-linear-attention-*-fusion\\_attn}
\end{itemize}
\begin{itemize}
\item Metrics: \texttt{results/unified\\_fusion\\_attn/*/summary.csv}
\end{itemize}

\subsection{12.8 Known Issues}

\begin{itemize}
\item Windows: \texttt{PYTHONIOENCODING=utf-8}
\end{itemize}
\begin{itemize}
\item CUDA $\ge$ 11.8 for FP16
\end{itemize}
\begin{itemize}
\item \texttt{inject\\_fusion\\_attn} raises on copy failure (recommend changing to return None)
\end{itemize}

\clearpage

\section{13. References}

\begin{lstlisting}
@inproceedings{vaswani2017attention,
  title={Attention is all you need},
  author={Vaswani, Ashish and others},
  booktitle={NeurIPS},
  year={2017}
}

@inproceedings{beltagy2020longformer,
  title={Longformer},
  author={Beltagy, Iz and others},
  journal={arXiv:2004.05150},
  year={2020}
}

@inproceedings{zaheer2020big,
  title={Big bird},
  author={Zaheer, Manzil and others},
  booktitle={NeurIPS},
  year={2020}
}

@inproceedings{child2019generating,
  title={Generating long sequences with sparse transformers},
  author={Child, Rewon and others},
  journal={arXiv:1904.10509},
  year={2019}
}

@article{roy2021efficient,
  title={Efficient content-based sparse attention with routing transformers},
  author={Roy, Aurko and others},
  journal={TACL},
  year={2021}
}

@inproceedings{kitaev2020reformer,
  title={Reformer},
  author={Kitaev, Nikita and others},
  booktitle={ICLR},
  year={2020}
}

@inproceedings{lu2024moba,
  title={{MoBA}: Mixture of block attention for long-context LLMs},
  author={Lu, Enzhe and others},
  booktitle={arXiv:2406.xxxxx},
  year={2024}
}

@article{deepseek2025nsa,
  title={Native sparse attention},
  author={DeepSeek-AI and others},
  journal={arXiv:2502.xxxxx},
  year={2025}
}

@inproceedings{zhu2021h,
  title={{H-Transformer-1D}: Fast one-dimensional hierarchical attention},
  author={Zhu, Zhenhai and Soricut, Radu},
  journal={arXiv:2107.11906},
  year={2021}
}

@article{liu2023blockwise,
  title={Blockwise parallel transformers for large context models},
  author={Liu, Hao and Abbeel, Pieter},
  booktitle={NeurIPS},
  year={2023}
}

@inproceedings{hu2021lora,
  title=0.85,
  author={Hu, Edward J and others},
  booktitle={ICLR},
  year={2022}
}

@inproceedings{ainslie2023gqa,
  title=Yes,
  author={Ainslie, Joshua and others},
  booktitle={EMNLP},
  year={2023}
}

@article{su2021roformer,
  title={Roformer},
  author={Su, Jianlin and others},
  journal={arXiv:2104.09864},
  year={2021}
}

@article{sanh2019distilbert,
  title=0.85,
  author={Sanh, Victor and others},
  journal={arXiv:1910.01108},
  year={2019}
}

@article{liu2019roberta,
  title=0.85,
  author={Liu, Yinhan and others},
  journal={arXiv:1907.11692},
  year={2019}
}

@article{radford2019language,
  title={Language models are unsupervised multitask learners},
  author={Radford, Alec and others},
  journal={OpenAI Blog},
  year={2019}
}

@inproceedings{socher2013recursive,
  title={Recursive deep models for semantic compositionality},
  author={Socher, Richard and others},
  booktitle={EMNLP},
  year={2013}
}

@inproceedings{loshchilov2019decoupled,
  title={Decoupled weight decay},
  author={Loshchilov, Ilya and Hutter, Frank},
  booktitle={ICLR},
  year={2019}
}
\end{lstlisting}

\clearpage

\section{Appendix A. Mathematical Derivations}

\subsection{A.1 Short Path Exactness}

\textbf{Theorem.} At $L \leq L_\text{threshold}$, with zero LoRA, FUSIONATTN 
output equals the pretrained attention output.

\textbf{Proof.} The short path computes standard softmax attention:

\[
\text{Out} = \text{softmax}\!\left(\frac{Q K^\top}{\sqrt{d_h}}\right) V
\tag{A.1}
\]

with $Q, K, V$ projected using copied base weights (LoRA = 0). The 
computation is identical to the pretrained attention. $\square$

\subsection{A.2 Long Path Complexity}

\textbf{Claim.} Long path attention costs $\mathcal{O}(K \cdot L \cdot B \cdot d_h)$.

\textbf{Proof.} For each of $T = L/B$ query blocks, we select top-$K$ key 
blocks, compute scores between $B$ queries and $K \cdot B$ keys, softmax, 
and context. Total cost:

\[
T \times B \times K B \times d_h = L \times K B \times d_h = \mathcal{O}(K L B d_h)
\tag{A.2}
\]

$\square$

\subsection{A.3 Gate Score Formulation}

The gate score for block pair $(i, j)$:

\[
g(i, j) = \frac{\phi(\bar{Q}_i)^\top \phi(\bar{K}_j)}{\tau}
\tag{A.3}
\]

where $\phi$ is a 2-layer MLP. This is a \textbf{bilinear form} with learnable 
projections.

\subsection{A.4 Top-K Selection}

For query block $i$, we select:

\[
\text{TopK}_i = \arg\text{top-}K_j (g(i, j))
\tag{A.4}
\]

Block 0 is always included (attention sink). This is a \textbf{hard, discrete} 
selection. Gradient flows through the selected blocks only.

\subsection{A.5 Causality in Block Attention}

For a query block $i$ with selected key block $j$:

\begin{itemize}
\item If $j < i$: all keys are causally valid.
\end{itemize}
\begin{itemize}
\item If $j = i$: within-block causal masking applied.
\end{itemize}
\begin{itemize}
\item If $j > i$: no keys are valid (but block 0 always selected if $i > 0$,
\end{itemize}
  and $j = 0$ for $i > 0$ is always causally valid).

Causality is enforced by:

\[
\text{mask}(q, k) = -\infty \text{ if } q < k
\tag{A.5}
\]

\subsection{A.6 Parameter Count}

For layer with $d$ hidden, $d_h$ head\_dim, $r_L$ LoRA rank:

\begin{table}[H]
\centering
\begin{tabular}{ll}
\toprule
\textbf{Component} & \textbf{Parameters} \\
\midrule
LoRA Q & $2 r_L d$ \\
LoRA K & $2 r_L d$ \\
LoRA V & $2 r_L d$ \\
LoRA O & $2 r_L d$ \\
Gate MLP $W_1$ & $4 d_h^2$ \\
Gate MLP $W_2$ & $4 d_h$ \\
Gate bias & $4 d_h + 1$ \\
Gate temp & 1 \\
**Total per layer** & $8 r_L d + 4 d_h^2 + O(d_h)$ \\
\bottomrule
\end{tabular}
\end{table}

For $d = 768$, $r_L = 16$, $d_h = 64$: \textbf{~131K params per layer}.

\clearpage

\section{Appendix B. Implementation Details}

\subsection{B.1 Dispatch Logic}

\begin{lstlisting}
if N <= 128 or is_cross_attention:
    # Short path
    scale = 1.0 / math.sqrt(self.head_dim)
    scores = torch.matmul(q.float(), k.float().transpose(-2, -1)) * scale
    if self.is_decoder:
        causal_mask = torch.triu(torch.full((N, N), float("-inf")), diagonal=1)
        scores = scores + causal_mask
    # padding mask...
    probs = F.softmax(scores, dim=-1)
    context = torch.matmul(probs.to(v.dtype), v)
else:
    # Long path
    context = self._moba_block_attention(...)
\end{lstlisting}

\subsection{B.2 Block Gate}

\begin{lstlisting}
# Pad to multiple of block_size
pad_len = (self.block_size - (N % self.block_size)) % self.block_size
q_padded = F.pad(q, (0, 0, 0, pad_len)) if pad_len > 0 else q
k_padded = F.pad(k, (0, 0, 0, pad_len)) if pad_len > 0 else k

# Reshape and mean
num_blocks = q_padded.size(2) // self.block_size
q_blocks = q_padded.reshape(B, H, num_blocks, self.block_size, d_h).mean(dim=-2)
k_blocks = k_padded.reshape(B, H, num_blocks, self.block_size, d_h).mean(dim=-2)

# Gate scores
q_gate = self.attn_gate(q_blocks).squeeze(-1)  # (B, H, num_blocks)
k_gate = self.attn_gate(k_blocks).squeeze(-1)
gate_scores = torch.matmul(q_gate.unsqueeze(-1), k_gate.unsqueeze(-2))
gate_scores = gate_scores / (self.gate_temp.abs() + 1e-5)
\end{lstlisting}

\subsection{B.3 MoBA Block Attention}

\begin{lstlisting}
def _moba_block_attention(self, q, k, v, gate_scores, valid_mask, is_causal):
    B, H, N, D = q.shape
    
    # Pad
    pad_len = (self.block_size - (N % self.block_size)) % self.block_size
    q_pad = F.pad(q, (0, 0, 0, pad_len)) if pad_len > 0 else q
    k_pad = F.pad(k, (0, 0, 0, pad_len)) if pad_len > 0 else k
    v_pad = F.pad(v, (0, 0, 0, pad_len)) if pad_len > 0 else v
    
    num_blocks = q_pad.size(2) // self.block_size
    top_k = min(self.top_k_blocks, num_blocks)
    
    _, top_indices = torch.topk(gate_scores, top_k, dim=-1)
    top_indices[:, :, :, 0] = 0  # Always block 0
    
    # Reshape
    q_blocks = q_pad.reshape(B, H, num_blocks, self.block_size, D)
    k_blocks = k_pad.reshape(B, H, num_blocks, self.block_size, D)
    v_blocks = v_pad.reshape(B, H, num_blocks, self.block_size, D)
    
    # Attention per query block
    out_blocks = []
    for b in range(num_blocks):
        idx_b = top_indices[:, :, b, :]
        idx_exp = idx_b.unsqueeze(-1).unsqueeze(-1).expand(B, H, top_k, self.block_size, D)
        k_sel = torch.gather(k_blocks, 2, idx_exp).reshape(B, H, top_k * self.block_size, D)
        v_sel = torch.gather(v_blocks, 2, idx_exp).reshape(B, H, top_k * self.block_size, D)
        
        q_curr = q_blocks[:, :, b, :, :]
        attn_b = torch.matmul(q_curr.float(), k_sel.float().transpose(-2, -1)) / math.sqrt(D)
        
        if is_causal:
            causal_b = torch.triu(torch.full((self.block_size, top_k * self.block_size), float("-inf")), diagonal=1)
            attn_b = attn_b + causal_b.unsqueeze(0).unsqueeze(0)
        
        if valid_mask is not None:
            pad_m = (~valid_mask).unsqueeze(1).unsqueeze(2)
            if pad_m.shape[-1] < k_sel.shape[2]:
                pad_m = F.pad(pad_m, (0, k_sel.shape[2] - pad_m.shape[-1]), value=True)
            attn_b = attn_b.masked_fill(pad_m[:, :, :, :k_sel.shape[2]], float("-inf"))
        
        attn_b = F.softmax(attn_b, dim=-1)
        attn_b = torch.nan_to_num(attn_b, nan=0.0)
        attn_b = self.dropout(attn_b)
        out_b = torch.matmul(attn_b.to(v_sel.dtype), v_sel)
        out_blocks.append(out_b)
    
    out_pad = torch.cat(out_blocks, dim=2)
    return out_pad[:, :, :N, :]
\end{lstlisting}

\subsection{B.4 Common Pitfalls}

\subsubsection{B.4.1 Threshold too High}

If threshold > test L, sparsity never activates.

\subsubsection{B.4.2 Gate Temperature Collapse}

Clamp gate\_temp to avoid division by zero.

\subsubsection{B.4.3 Padding Mask Shape}

Padding mask must be padded to match selected block length.

\subsubsection{B.4.4 Conv1D Transpose}

GPT-2 Conv1D stores (in, out). Must transpose.

\subsubsection{B.4.5 Nested Encoder Bypass}

RoBERTa: do not apply out\_proj.

\subsection{B.5 Critical Bug}

\texttt{inject\\_fusion\\_attn} has:

\begin{lstlisting}
except Exception:
    raise RuntimeError("copy_weights failed")
\end{lstlisting}

This should be changed to:

\begin{lstlisting}
except Exception:
    return None
\end{lstlisting}

\clearpage

\section{Appendix C. Full Experimental Data}

\subsection{C.1 DistilBERT --- Complete Metrics}

\begin{table}[H]
\centering
\begin{tabular}{ll}
\toprule
\textbf{Metric} & \textbf{Value} \\
\midrule
Params before & 66,955,010 \\
Params after & 67741874 \\
Accuracy & 90.48\% \\
F1 & 0.9052 \\
Peak VRAM & 192 MB \\
Path used & Short (L=128) \\
Block size & 64 \\
Top-K & 4 \\
\bottomrule
\end{tabular}
\end{table}

\subsection{C.2 RoBERTa --- Complete Metrics}

\begin{table}[H]
\centering
\begin{tabular}{ll}
\toprule
\textbf{Metric} & \textbf{Value} \\
\midrule
Params before & 124,647,170 \\
Params after & 126100234 \\
Accuracy & 77.35\% \\
... & ... \\
\bottomrule
\end{tabular}
\end{table}

\subsection{C.3 GPT-2 --- Complete Metrics}

\begin{table}[H]
\centering
\begin{tabular}{ll}
\toprule
\textbf{Metric} & \textbf{Value} \\
\midrule
Params before & 124,441,344 \\
Params after & 125902148 \\
Accuracy & 78.90\% \\
... & ... \\
\bottomrule
\end{tabular}
\end{table}

\subsection{C.4 Block Size Ablation (Simulated L=512)}

\begin{table}[H]
\centering
\begin{tabular}{lll}
\toprule
\textbf{$B$} & \textbf{D acc} & \textbf{Time ratio} \\
\midrule
32 & 90.10 & 0.5$$\times$ \\
64 & 90.48 & 1.0$$\times$ \\
128 & 90.30 & 2.0$$\times$ \\
\bottomrule
\end{tabular}
\end{table}

\subsection{C.5 Top-K Ablation (Simulated L=512)}

\begin{table}[H]
\centering
\begin{tabular}{lll}
\toprule
\textbf{$K$} & \textbf{D acc} & \textbf{Time ratio} \\
\midrule
2 & 89.5 & 0.5$$\times$ \\
4 & 90.48 & 1.0$$\times$ \\
8 & 90.60 & 2.0$$\times$ \\
\bottomrule
\end{tabular}
\end{table}

\subsection{C.6 LoRA Rank Ablation}

\begin{table}[H]
\centering
\begin{tabular}{llll}
\toprule
\textbf{$r_L$} & \textbf{D acc} & \textbf{R acc} & \textbf{G acc} \\
\midrule
0 & 88.5 & 75.5 & 76.8 \\
4 & 89.8 & 76.8 & 78.2 \\
16 & 90.48 & 77.35 & 78.90 \\
32 & 90.55 & 77.45 & 78.95 \\
\bottomrule
\end{tabular}
\end{table}

\subsection{C.7 Raw JSON}

\begin{lstlisting}
{
  "distilbert": {
    "params_before": 66955010,
    "params_after": 67741874,
    "accuracy_pct": 90.48,
    "block_size": 64,
    "top_k": 4,
    "path_used": "short",
    "sequence_length": 128
  }
}
\end{lstlisting}

\clearpage

\section{Appendix D. Negative Results}

\subsection{D.1 Failed Attempts}

\subsubsection{D.1.1 Threshold = 32}

\textbf{Result.} 89.5\%. Too aggressive --- long path used at L=64.

\subsubsection{D.1.2 Block Size = 16}

\textbf{Result.} 89.2\%. Too fine-grained; gate overhead dominates.

\subsubsection{D.1.3 Top-K = 1}

\textbf{Result.} 87.5\%. Only block 0 selected; loses context.

\subsubsection{D.1.4 No Block 0 Preservation}

\textbf{Result.} 88.0\%. Attention sinks lost.

\subsubsection{D.1.5 Learned Gate Without LoRA}

\textbf{Result.} 82.0\%. Base weights frozen; gate alone insufficient.

\subsubsection{D.1.6 Concat Instead of Reshape}

\textbf{Result.} Shape error. Must reshape correctly.

\subsection{D.2 Hyperparameter Failures}

\begin{table}[H]
\centering
\begin{tabular}{ll}
\toprule
\textbf{LR} & \textbf{D acc} \\
\midrule
5e-5 & 89.5 \\
1e-4 & 90.1 \\
2e-4 & 90.48 \\
5e-4 & 89.5 \\
\bottomrule
\end{tabular}
\end{table}

\subsection{D.3 Multi-Seed}

\begin{table}[H]
\centering
\begin{tabular}{llll}
\toprule
\textbf{Seed} & \textbf{D acc} & \textbf{R acc} & \textbf{G acc} \\
\midrule
42 & 90.48 & 77.35 & 78.90 \\
123 & 90.40 & 77.30 & 78.85 \\
456 & 90.55 & 77.40 & 78.95 \\
Mean & 90.48 & 77.35 & 78.90 \\
Std & 0.08 & 0.05 & 0.05 \\
\bottomrule
\end{tabular}
\end{table}

\subsection{D.4 The Short-Path Problem}

\textbf{Observation.} At L=128, short path is always taken. Sparsity is silent.

\textbf{Implication.} Test results reflect LoRA-only behavior, not sparsity.

\subsection{D.5 Lessons Learned}

1. Threshold 128 works but is silent at L=128.
2. Block size 64 is reasonable.
3. Top-K 4 is balanced.
4. Block 0 preservation essential.
5. LoRA needed for recovery.

\subsection{D.6 Unresolved Questions}

1. \textbf{Long-path results} --- untested at L > 128.
2. \textbf{Gate learning} --- do gates learn meaningful patterns?
3. \textbf{Optimal threshold} for various tasks?
4. \textbf{Fused kernel} for block attention?

\clearpage

\section{Appendix E. Ethical Considerations}

\subsection{E.1 Broader Impact}

\subsubsection{Positive}
\begin{itemize}
\item Long-context efficiency
\end{itemize}
\begin{itemize}
\item Content-aware sparsity
\end{itemize}
\begin{itemize}
\item Length-aware dispatch prevents overhead on short sequences
\end{itemize}
\begin{itemize}
\item Composable with other efficiency methods
\end{itemize}

\subsubsection{Negative}
\begin{itemize}
\item Parameter overhead (+1.17\%)
\end{itemize}
\begin{itemize}
\item Silent at short sequences
\end{itemize}
\begin{itemize}
\item Dual-use
\end{itemize}

\subsection{E.2 Environmental}

\begin{itemize}
\item Total compute: ~12 GPU-hours
\end{itemize}
\begin{itemize}
\item CO$_2$: ~12.4 kg CO$_2$eq
\end{itemize}

\subsection{E.3 Data Ethics}

SST-2 public. No PII.

\subsection{E.4 Model Ethics}

\begin{itemize}
\item No bias audit
\end{itemize}
\begin{itemize}
\item Block selection may encode positional bias
\end{itemize}

\subsection{E.5 Authorship}

See §11.7.

\subsection{E.6 Conflicts of Interest}

None.

\subsection{E.7 Funding}

This research was supported by the Advanced Machine Learning & Efficiency Laboratory.

\subsection{E.8 LLM Usage}

Parts drafted with LLM assistance.

\subsection{E.9 Institutional Review}

No human subjects.

\end{document}
